% Patrick Bijsmans: Euroskepticism, a Multifaceted Phenomenon – Reading Summary
% Introduction to the European Union, Week 5
% Compile with: pdflatex bijsmans-euroskepticism-reading-summary.tex

\documentclass[11pt,a4paper]{article}
\usepackage[utf8]{inputenc}
\usepackage[T1]{fontenc}
\usepackage[margin=2.5cm]{geometry}
\usepackage{booktabs}
\usepackage{enumitem}
\setlist{nosep,leftmargin=*}
\usepackage{parskip}
\usepackage{microtype}
\usepackage{hyperref}
\usepackage{xcolor}
\definecolor{eublue}{RGB}{0,51,153}

\hyphenation{Spitzenkandidat Spitzenkandidaten}
\hypersetup{
  colorlinks=true,
  linkcolor=eublue,
  urlcolor=eublue,
  pdftitle={Bijsmans – Euroskepticism},
  pdfauthor={Introduction to the EU, UCM}
}

\begin{document}

\title{Patrick Bijsmans: Euroskepticism, a Multifaceted Phenomenon}
\author{Introduction to the European Union\\Universidad Complutense de Madrid\\Week 5}
\date{}
\maketitle
\vspace{-1em}
\hrule
\vspace{1em}

\section*{Rhetorical Framework (Ethos, Logos, Pathos)}

\begin{table}[h]
\centering
\begin{tabular}{@{}lll@{}}
\toprule
\textbf{Appeal} & \textbf{Meaning} & \textbf{Bijsmans's Use} \\
\midrule
\textbf{Ethos} & Credibility, authority, trust & Encyclopedic synthesis; integrates parties, public opinion, crises; avoids simplistic labels \\
\textbf{Logos} & Logic, structure, argument & Conceptual map (soft/hard, levels, dimensions); causal factors (domestic politics, crises) \\
\textbf{Pathos} & Emotion, stakes, persuasion & Implicit---Euroskepticism as politically charged; appeal to readers seeking nuance over polemic \\
\bottomrule
\end{tabular}
\end{table}

\textbf{Key insight:} Bijsmans prioritizes \textit{logos} (conceptual clarity, multidimensional framework) to build \textit{ethos} (authoritative, balanced account). \textit{Pathos} is mostly implicit---the stakes of Euroskepticism (legitimacy, crisis, populism) are present but not dramatized.

\section*{Bibliographic Information}

\begin{itemize}
\item \textbf{Author:} Patrick Bijsmans
\item \textbf{Title:} Euroskepticism, a multifaceted phenomenon
\item \textbf{Source:} Oxford Research Encyclopedia of Politics, Oxford University Press
\item \textbf{Year:} 2020
\end{itemize}

\section*{Summary \& Key Points}

\begin{itemize}
\item Euroskepticism is not a single attitude but a spectrum of positions, from mild criticism of specific EU policies to principled rejection of integration.
\item The article distinguishes between different types (e.g.\ soft vs.\ hard) and levels (parties, elites, citizens) of Euroskepticism.
\item Domestic politics, socio-economic conditions and EU-level crises (euro, migration, rule-of-law disputes) all shape how and when Euroskepticism becomes politically powerful.
\end{itemize}

\section*{Main Arguments}

\subsection*{I. Euroskepticism as Multidimensional}

Euroskepticism should be understood as a multidimensional phenomenon, not just a simple ``for or against'' the EU. Different actors may oppose different aspects of integration (e.g.\ economic policies vs.\ political institutions vs.\ identity) and for different reasons.

\subsection*{II. Domestic Politics and Media Framing}

Bijsmans argues that party competition and media framing in domestic arenas are central to explaining how Euroskepticism develops and spreads. National political contexts mediate citizens' views on the EU.

\section*{Relevance to Course}

\begin{itemize}
\item \textbf{Ethos:} Provides vocabulary and conceptual tools for analyzing contemporary debates.
\item \textbf{Logos:} Soft/hard, levels, dimensions---framework for interpreting variation across member states.
\item \textbf{Pathos:} Post-2008 crises, Brexit, rule-of-law conflicts---stakes of Euroskepticism.
\item \textbf{All three:} Helps interpret party positions, referendums, election outcomes; connects to legitimacy and democratic contestation.
\end{itemize}

\section*{Discussion Questions}

\begin{itemize}
\item Euroskepticism as democratic corrective or destabilizing force? (Ethos: who frames it? Logos: evidence, causal claims; Pathos: stakes of legitimacy)
\item How do domestic party systems shape Euroskepticism? (Ethos: national contexts; Logos: causal mechanisms; Pathos: political competition)
\item Can transparency or participation reduce Euroskepticism? (Ethos: who proposes reforms? Logos: what evidence? Pathos: hope, skepticism)
\end{itemize}

\vspace{1em}
\noindent\footnotesize\textit{Reference:} Bijsmans, Patrick (2020). ``Euroskepticism, a multifaceted phenomenon.'' \textit{Oxford Research Encyclopedia of Politics}, Oxford University Press.

\end{document}
